\section{Fleet Learning Dynamics}

\label{sec:convergence}

\subsection{Experience Quality Metric}

\begin{definition}[Fleet Experience Quality]
\label{def:quality-metric}
The fleet-wide experience quality at round $T$ is:
\begin{equation}
    Q^{(T)} = \frac{1}{N} \sum_{i=1}^{N} \frac{1}{|\mathcal{L}_i^{(T)}|}
    \sum_{E_m \in \mathcal{L}_i^{(T)}} q_m \cdot \Sim(E_m, \mathcal{T}_i),
    \label{eq:fleet-quality}
\end{equation}
where $\Sim(E_m, \mathcal{T}_i)$ is the cosine similarity between the
experience embedding $\bm{e}_m$ and the mean embedding of node $i$'s
recent task distribution $\mathcal{T}_i$.
\end{definition}

$Q^{(T)}$ captures both the intrinsic quality of experiences (via $q_m$) and
their relevance to each node's task distribution (via $\Sim$). A fleet where
nodes only receive experiences from irrelevant domains would have high $q_m$
but low $\Sim$, correctly receiving a low $Q$ score.

\begin{theorem}[Monotonic Fleet Improvement]
\label{thm:monotone}
Under Assumptions~\ref{ass:threat}--\ref{ass:diversity}, if DeltaSoup's
filter threshold $\tau_{\text{filter}}$ is set such that Byzantine experiences
are excluded with probability $\geq 1 - \delta$ (as in
Theorem~\ref{thm:deltasoup-robust}), then the expected fleet quality is
non-decreasing:
\begin{equation}
    \mathbb{E}\left[Q^{(T+1)}\right] \geq \mathbb{E}\left[Q^{(T)}\right]
    - O(\delta).
    \label{eq:monotone}
\end{equation}
\end{theorem}

\begin{proof}
At round $T+1$, each node receives honest experiences from its peers and
merges them via DeltaSoup. Any honest experience $E_m$ with $\Sim(E_m,
\mathcal{T}_i) > 0$ and $q_m > 0$ increases the quality of node $i$'s
library, contributing positively to $Q^{(T+1)}$. Byzantine experiences that
pass the filter (probability $\leq \delta$ per experience) may decrease
quality. Taking expectations and applying linearity, the expected quality
change is $\geq -O(\delta) \cdot |\mathcal{E}_{\text{Byzantine}}|$. Since
$|\mathcal{E}_{\text{Byzantine}}| = O(f)$ and $f < N/3$, the net expected
change is non-negative for sufficiently small $\delta$.
\end{proof}

\subsection{Information Flow Model}

The fleet learning cycle for a single round is:
\begin{equation}
    \underbrace{\text{Agent Session}}_{\text{tool calls, outcomes}}
    \xrightarrow{\text{GRPO}}
    \underbrace{\text{Local Experiences}}_{\mathcal{L}_i^{(T)}}
    \xrightarrow{\text{BitDelta}}
    \underbrace{\text{Compressed Batch}}_{\hat{B}_i}
    \xrightarrow{\text{Gossip}}
    \underbrace{\text{Peer Receives}}_{\hat{B}_j, j \neq i}
    \xrightarrow{\text{DeltaSoup}}
    \underbrace{\text{Updated Library}}_{\mathcal{L}_j^{(T+1)}}
    \label{eq:info-flow}
\end{equation}

\paragraph{Compression-Decompression Semantic Preservation.}
The information flow~\eqref{eq:info-flow} introduces two lossy steps:
(i) BitDelta compression (Theorem~\ref{thm:reconstruction}) and
(ii) DeltaSoup filtering and merging. By Corollary~\ref{cor:cosine},
BitDelta preserves cosine similarity $\geq 0.95$. The DeltaSoup merge
(Eq.~\ref{eq:weighted-merge}) is a convex combination within the cluster,
so the merged embedding remains close to all cluster members. Combined,
the end-to-end semantic preservation satisfies:
\begin{equation}
    \cos(\bm{e}_m, \bm{e}_m^{\text{final}}) \geq 0.95 \cdot 0.95 = 0.9025,
    \label{eq:end-to-end-cosine}
\end{equation}
where $\bm{e}_m^{\text{final}}$ is the embedding after compression,
broadcast, decompression, and DeltaSoup merge.

\subsection{Specialization vs.\ Generalization}

\paragraph{Domain-Aware Routing.}
Each experience carries a domain tag (e.g., \texttt{coding}, \texttt{writing},
\texttt{data-analysis}). During broadcast, nodes may optionally subscribe to
specific domain tags, filtering incoming batches to relevant domains only.
This provides two benefits:
\begin{enumerate}[leftmargin=1.5em]
    \item \textbf{Bandwidth reduction:} A coding agent need not receive
          marketing-domain experiences.
    \item \textbf{Quality improvement:} Domain-filtered libraries have higher
          mean $\Sim(E_m, \mathcal{T}_i)$, directly improving $Q^{(T)}$.
\end{enumerate}

\paragraph{Fleet Diversity.}
Despite domain specialization, the fleet benefits from cross-domain transfer
when tasks span multiple domains. DeltaSoup's confidence-weighted merge
naturally weights domain-relevant experiences more heavily (as relevant
experiences tend to have higher $q_m$ from the receiving node's evaluations).

\begin{proposition}[Specialization-Generalization Tradeoff]
\label{prop:tradeoff}
Let $\rho \in [0,1]$ be the fraction of received experiences that are in the
node's primary domain. The fleet quality satisfies:
\begin{equation}
    Q^{(T)}(\rho) = \rho \cdot Q_{\text{specialized}}^{(T)}
                  + (1 - \rho) \cdot Q_{\text{general}}^{(T)},
\end{equation}
where $Q_{\text{specialized}}^{(T)} \geq Q_{\text{general}}^{(T)}$ for
single-domain tasks and $Q_{\text{specialized}}^{(T)} \leq Q_{\text{general}}^{(T)}$
for multi-domain tasks. The optimal $\rho^*$ depends on task distribution breadth.
\end{proposition}

% ============================================================================
% 7. INTEGRATION WITH AGENT HARNESS
% ============================================================================
